\section{The small-noise half-amplitude step without more cells}
\label{sn:section}

Fix an actual integer $n\ge42$. Let $(P,H)$ be a fixed-target maximizing
pair with nonnegative excess $\mathfrak a$ over $c=C_E$, and assume the
true predecessor bound $C_{n-1}\le c$. Retain the closed-support contact
theorem. Suppose the already acquired panel bounds give
\[
 p=\Pp(X>0)\in[41/100,42/100],\quad
 X=\ell(Y-p),\quad Y=B+U,\quad B=\mathbf1_{X>0},\quad
 \E(U\mid B)=0,
\]
\[
 |U|<3/5,\qquad U_0<d<69/1000,\qquad U_1>-d,
 \qquad\lambda=n\E U^2\le3/40.
\]
Here $X$ is centered with variance one and $\ell>0$. In particular the
initial strip is established before this argument uses bounded conditional
moments. These are statements about one actual marginal; no relaxation is
assumed to be a probability law.

\begin{theorem}\label{sn:theorem}
Under these hypotheses $|U|<1/2$. The conclusion holds for every raw
threshold $H\in\mathbb R$, both discrepancy orientations after their
proper closed reflection, and every actual index $n\ge42$. It includes
$\lambda=0$ without taking a positive-noise limit.
\end{theorem}

\subsection{Moment bounds that vanish with the actual noise}
Write $q=1-p$, $v=pq$, $s=\E U^2=\lambda/n$ and
$\theta=\E(X-\E[X\mid B])^2$. Conditional centering and variance give
\[
 \ell=(v+s)^{-1/2},\quad
 \theta=\frac{s}{v+s}=\frac{\lambda}{nv+\lambda},\quad
 m^2=v(1-\theta),\quad m=\E X_+.
\]
Put $v_-=2419/10000$, $v_+=609/2500$ and
$\Theta=733/100000$. Then
\[
 \theta\le\frac{3/40}{42v_-+3/40}<\Theta,
 \qquad \ell>\frac{1009}{500}=2.018,
\]
where the latter follows by the exact square comparison
$(1009/500)^2v_+<1-\Theta$.

The sign of $B-p+U$ is the actual sign used to define $B$. Expanding
the absolute cubic on these two actual sign classes, and using conditional
centering, yields the exact identity
\[
 \rho=\ell^3\left\{v(1-2v)+3\E[|B-p|U^2]
                  +\E[\operatorname{sgn}(B-p)U^3]\right\}.
\]
Since $|U|\le3/5$ and $p<1/2$,
\[
 \rho\le\rho_B(p)+\frac{3(1-p)+3/5}{v^{3/2}}\frac\lambda n,
 \qquad \rho_B(p)=\frac{1-2pq}{\sqrt{pq}}.
\]
The function $\rho_B$ decreases on this panel. The exact comparisons
\[
 24^2v_-^3>\{3(1-41/100)+3/5\}^2,
 \qquad (21/20)^2v_->(1-2v_-)^2
\]
therefore imply
\begin{equation}
 \rho<\frac{21}{20}+24\frac{3/40}{42}
       =\frac{21}{20}+\frac3{70}<\frac{1093}{1000}=:r.
 \label{sn:rho}
\end{equation}
This proves the coefficient24 from the previously acquired strip; it
does not assume the half-amplitude conclusion.

Let $\theta_\pm$ be the two unnormalized conditional variance contributions.
They are nonnegative and sum to $\theta$. Consequently
\[
 \beta=\E[X|X|]=1-2p-2q\theta+2\theta_+,
 \quad 1-2p-2q\theta\le\beta\le1-2p+2p\theta.
\]
Both bounding functions, with $\theta$ replaced outward by $\Theta$,
decrease in $p$. Thus
\begin{equation}
 \frac{151}{1000}<1-2(42/100)-2(58/100)\Theta
 \le\beta\le1-2(41/100)+2(41/100)\Theta<\frac{187}{1000}.
 \label{sn:beta}
\end{equation}
No assertion of uniform integrability, optimizer structure, or independent
conditional noise laws is used in these single-marginal moment identities.

\subsection{Uniform Gaussian contact bounds at the new boundaries}
At every point $t$ of the closed support, fixed-target contact gives
\begin{equation}
 0\le\mathcal H_n(H,t)+c\{-|t|^3+\tfrac32\rho t^2+3\beta t+\rho e_n\}
                  -n\mathfrak a,\qquad
 e_n=n(\sqrt{n/(n-1)}-1)-\tfrac12.
 \label{sn:contact}
\end{equation}
Set $v_n=n-1$, $s_n=\sqrt{n/v_n}$ and $F_v(x)=\Phi(x/\sqrt v)$.
The exact heat expression is
\[
 \frac{\mathcal H_n(H,t)}{n^{3/2}}
 =F_{v_n}(H-t)-F_n(H)+tF_n'(H)-\frac{t^2-1}{2}F_n''(H).
\]
Expand the spatial term through degree four. In the other three terms,
expand in the variance variable to orders two, one, one, using
$\partial_vF_v=F_v''/2$. After cancellation,
\[
 \frac{\mathcal H_n(H,\pm u)}{s_n^3}
 =\pm\frac{u^3-3u}{6}(1-z^2)\phi(z)
  +\frac{u^4-6u^2+3}{24\sqrt{v_n}}\phi'''(z)+R,
 \quad z=H/\sqrt{v_n}.
\]
The four Taylor remainders are bounded by
\[
 \frac{u^5\|\phi^{(4)}\|_\infty}{120v_n},\quad
 \frac{u\|\phi^{(4)}\|_\infty}{8v_n},\quad
 \frac{\|\phi^{(5)}\|_\infty}{48v_n^{3/2}},\quad
 \frac{|u^2-1|\|\phi^{(5)}\|_\infty}{16v_n^{3/2}}.
\]
Fourier inversion of the Gaussian gives
$\|\phi''\|_\infty\le\phi(0)$,
$\|\phi'''\|_\infty\le2/\pi<2\phi(0)$,
$\|\phi^{(4)}\|_\infty\le3\phi(0)<6/5$, and
$\|\phi^{(5)}\|_\infty\le8/\pi<51/20$.
These follow by integrating the absolute powers of frequency; the last
comparison uses $\pi>160/51$. Thus for every $u\ge1$,
\[
 |R|\le\frac{u^5/100+3u/20}{v_n}
          +\frac{51(3u^2-2)}{960v_n^{3/2}}.
\]
This is the precise domain extension: the remainder requires only
$u\ge1$, and the positive leading coefficient requires $u>\sqrt3$.
The sign of $u^4-6u^2+3$ is handled by its absolute value.

For $n\ge42$, $s_n<81/80=:S$, $\sqrt{v_n}>6$ and $e_n<8/875$.
For the last bound, $s_n<101213/100000$ by squaring, and the increasing
function $s^2/(s+1)-1/2$ is less than $8/875$ at that rational upper.
Set $c_-=819/2000<c$ and
\begin{align*}
 E(u)&=u^5/4100+3u/820+17(3u^2-2)/78720,\\
 Q_{r,b}(u)&=-u^3+\tfrac32ru^2+3bu+(8/875)r,\\
 P_{r,b}(u)&=\tfrac25S^3\left\{\frac{u^3-3u}{6}
                     +\frac{|u^4-6u^2+3|}{72}\right\}
                    +S^3E(u)+c_-Q_{r,b}(u).
\end{align*}
At a boundary with $Q_{r,b}<0$, this is an upper bound for the right
side of \eqref{sn:contact}, after dropping $-n\mathfrak a\le0$.
For the negative ray use $b=-151/1000$, and for the positive ray
$b=187/1000$, by \eqref{sn:beta}. The multiplier substitution is made
only after checking the sign of $Q$.

Here the two rational boundaries are
\[
 u_-=459/250=1.836,\qquad u_+=2179/1000=2.179.
\]
Both exceed $\sqrt3$. Independent exact arithmetic verifies
\[
 P_{r,-151/1000}(u_-)<0,\qquad P_{r,187/1000}(u_+)<0.
\]
For full rays, the derivative is bounded uniformly in H by
\[
 |\partial_t\mathcal H_n(H,t)|
 \le\phi(0)\{t^2/2+1/2+e_n\}.
\]
Indeed separate the second-order spatial remainder of the normal density
at variance n, bounded by $\phi(0)t^2/(2n^{3/2})$, from its variance
change. The heat integral for the latter, multiplied by $n^{3/2}$,
is at most $\phi(0)n(s_n-1)=\phi(0)(1/2+e_n)$.
At each boundary the checker verifies the six strict signs
\[
 Q,Q',Q'',L,L',L''<0,\qquad
 L(u)=\tfrac25(u^2/2+1/2+8/875)+c_-Q'(u).
\]
Since $Q'''=-6$ and $L''=2/5-6c_-<0$, these signs persist forever.
With the moment caps fixed, the contact expression decreases on each
ray from its strictly negative boundary. Hence the entire closed rays
$(-\infty,-u_-]$ and $[u_+,\infty)$ are absent from the support.

\subsection{Amplitude and the closed seam}
Finally,
\[
 u_-<\frac{1009}{500}(41/100+1/2),\qquad
 u_+<\frac{1009}{500}(3/2-42/100).
\]
Therefore $-1/2<Y<3/2$ throughout the support. Together with the
previously acquired inward bounds $U_0<d<0.069$ and $U_1>-d$, this
gives $|U|<1/2$. If $\lambda=0$, conditional variances are zero and
$U=0$ almost surely directly; the proof never divides by $\lambda$.

To join this lemma to the existing eight equal noise bands, use it on
the entire closed band $[0,3/40]$, for every H. Retain the separately
verified contact/resource gates on the other seven closed bands covering
$[3/40,3/5]$. Their shared endpoint $3/40$ is included in both covers.
The lemma does not assert those other gates pass; their complete source
binding and finite inequalities remain separate premises. No added noise
or phase cell is required for this alternative composition.
